% ─────────────────────────────────────────────────────────────────────
% Evaluation subsection (input before \S discussion). Numbers from the
% 2026-09-04 scaled run; see benchmarks/README.md for the raw tables.
% ─────────────────────────────────────────────────────────────────────
\subsection{Escalation Benchmark}
\label{sec:eval-escalation}

\paragraph{Setup.}
We compare routing strategies on the same engine, tools, and prompts; only the escalation strategy differs.
Workloads: sanitized MBPP~\cite{austin2021mbpp}, 200 problems beyond the first 30 used during development, and all 164 HumanEval problems~\cite{chen2021humaneval}.
Each problem provides a sandbox with its tests; the agent writes a solution file and may run the tests; success is decided by the harness re-running the tests in a fresh process, never by the agent's own claim.
Small model: Llama~3.1~8B; large model: Gemini~2.5~Flash.
As robustness checks we repeat the comparison with GPT-4o-mini as the large model and run a 20-problem spot check with GPT-4o.
The small-model-only counterfactual is run three times per problem to estimate $p_{\text{SLM}}$, the probability that the small model solves the problem unaided; Joule and the FrugalGPT-style cascade are run three times and reported as mean $\pm$ standard deviation.

\paragraph{Baselines.}
\textbf{SLM-only} and \textbf{LLM-only} bound the space.
\textbf{Naive cascade} reruns the whole task on the large model when the small-model run does not complete (EcoAssistant-style~\cite{zhang2023ecoassistant}).
\textbf{FrugalGPT-style cascade}~\cite{chen2023frugalgpt} scores the small-model answer with a cheap scorer and reruns on the large model below a threshold.
\textbf{AutoMix-style}~\cite{madaan2023automix} has the small model verify its own answer three times and escalates on a majority ``no''.
\textbf{Pre-router} picks a tier per task with one cheap classification call (RouteLLM-style~\cite{ong2024routellm}).
\textbf{Static router} is Joule's previous plan-then-execute pipeline (\S\ref{sec:routing}).
Scorer, verifier, and router calls are charged to the strategy that makes them.

\paragraph{Metrics.}
Success rate; cost per task relative to LLM-only; and routing quality against the counterfactual labels:
a task \emph{needs} escalation if $p_{\text{SLM}} < 0.5$ and the large model succeeds (81 of 200 problems);
a strategy \emph{escalates} on a task if it spends any large-model tokens on it.
Precision and recall follow; \emph{wasted} counts escalations on tasks with $p_{\text{SLM}} \geq 0.8$.
For Joule we also report consultation success (consultations after which the small model finished without a handoff) and handoff success.

\paragraph{Results.}
Table~\ref{tab:escalation-mbpp} reports the main pair.
The small model alone solves 55\% (mean of three runs); the large model alone 97\%.
Joule reaches 96\% ($\pm$1.0 over three seeds) at 39\% of large-model cost, matching the best cascade within noise while the cascades spend 53--66\%, with the highest recall of the escalating strategies (85\% $\pm$ 5) at comparable precision (68\% $\pm$ 1.2).
Of the problems where Joule involved the large model, 89\% of handoffs succeeded and 40\% of consultations let the small model finish without a handoff.
The previous plan-then-execute pipeline and the RouteLLM-style pre-router both stay near the small model's own success rate: a router that predicts difficulty before seeing any execution evidence adds almost nothing.
With GPT-4o-mini as the large model (84\% alone on this slice), Joule is the cheapest strategy at 46\% of large-model cost and 83\% success, one point below the large model alone and three below the FrugalGPT-style cascade, which spends 48\% more; the weaker large model lowers handoff success to 36\%, so the quality of the upper rung matters.
On the 20-problem spot check with GPT-4o, Joule matched its 95\% success at 27\% of its cost.

\begin{table}[t]
\centering
\small
\caption{MBPP, 200 unseen problems, Llama~3.1~8B $\rightarrow$ Gemini~2.5~Flash. Cost is relative to LLM-only. Precision and recall are measured against $p_{\text{SLM}}$ labels from three SLM-only runs per problem; Joule and FrugalGPT-style are means over three seeds.}
\label{tab:escalation-mbpp}
\begin{tabular}{lrrrrrr}
\toprule
Strategy & Success & Cost/LLM & LLM used & Precision & Recall & Wasted \\
\midrule
SLM-only & 55\% & 0.07 & 0\% & -- & -- & -- \\
LLM-only & 97\% & 1.00 & 100\% & -- & -- & -- \\
Static router & 57\% & 0.14 & 0\% & -- & 0\% & 0 \\
Naive cascade & 95\% & 0.53 & 42\% & 76\% & 78\% & 6 \\
FrugalGPT-style & 96\% & 0.60 & 50\% & 66\% & 81\% & 12 \\
AutoMix-style & 97\% & 0.66 & 55\% & 65\% & 88\% & 17 \\
Pre-router & 55\% & 0.07 & 0\% & -- & 0\% & 0 \\
\textbf{Joule (adaptive)} & \textbf{96\%} & \textbf{0.39} & 51\% & 68\% & 85\% & 14 \\
\bottomrule
\end{tabular}
\end{table}

\paragraph{HumanEval.}
On all 164 HumanEval problems~\cite{chen2021humaneval} with the same pair, the small model alone solves 34\% and the large model 96\%.
Joule matches the large model at 96\% for 64\% of its cost (recall 84\%, precision 72\%); the FrugalGPT-style cascade reaches 93\% at 72\%.
With a small model this weak, handoffs carry the result (94\% succeed) and consultations rarely suffice (22\%): consultation pays when the small model is close to a solution, handoff when it is not.

\paragraph{Long-horizon tasks.}
To test the claim that trajectory-level escalation separates from request-level cascades as tasks lengthen, we bundle four MBPP problems into one task: implement four functions in one module, all tests must pass.
Trajectories run 8--15 steps on average (up to 40) instead of 3--6.
On 30 such tasks with the main pair, Joule reaches 67\% at 47\% of the large model's cost; the FrugalGPT-style cascade reaches 57\% at 76\%, and the large model alone 53\%.
The small model's incremental work is what makes the difference: a handoff carries a partly built module with passing tests, so the large model finishes rather than restarts, whereas a cascade reruns the whole task.
Consultations rarely unblock the small model on these tasks (2 of 24) and handoffs do (12 of 22).
Three policy rules that short tasks never exercise were needed: a rising test pass fraction counts as progress; a test run that repeats a verified result is not a second failure; and failures are counted within a window of recent steps rather than over the whole run.

\paragraph{Three rungs.}
On 50 further MBPP problems with Gemini~2.5~Flash as a middle rung and GPT-4o as the top, the ladder policy reaches 94\% at 15\% of GPT-4o's cost, against 86\% at 39\% for the two-tier policy that hands the small model straight to GPT-4o, and 90\% for GPT-4o alone.
The ladder used the middle rung on 18 problems and the top rung on 6; 80\% of its handoffs succeeded, against 46\% for direct handoffs to the frontier model.
Flash alone solves 96\% of this slice for 13\% of GPT-4o's cost, so the cheapest correct choice here is the middle model by itself; the ladder's contribution is reaching that outcome without knowing in advance which rung suffices.

\paragraph{Ablations.}
Table~\ref{tab:ablations} removes one design choice at a time on 50 further MBPP problems with the main pair (three SLM-only runs per problem for the labels).
Without consultations, success is unchanged but cost rises from 24\% to 40\% of the large model's and the large model touches 88\% of tasks instead of 38\%: every escalation is a handoff and most are unnecessary (precision 30\%).
Without deterministic verification the policy sees almost no failures, escalates on 6\% of tasks, and 8 of 50 runs finish ``completed'' with wrong code; success falls to 80\%.
With the model's self-reported confidence in place of the evidence-based composite, the model claimed a confidence of at least 0.8 in 72 of the 107 decisions that immediately followed a failed verification; the hard triggers still fire, so success only dips to 96\%, but escalations rise (46\% of tasks) and cost by 17\%.
Prose-only consultations (the earlier design) keep success at 98\% but only 40\% of them let the small model finish, against 83\% when the consultant returns the edit; handoffs fall from 12 to 4 and cost from 42\% to 24\%.
Static checks are a small effect on three-line functions (96\% versus 98\%).

\begin{table}[t]
\centering
\small
\caption{Ablations on 50 MBPP problems, Llama~3.1~8B $\rightarrow$ Gemini~2.5~Flash. Each row is the full policy with one choice removed. Consult ok: consultations after which the small model finished without a handoff.}
\label{tab:ablations}
\begin{tabular}{lrrrrrr}
\toprule
Variant & Success & Cost/LLM & LLM used & Precision & Recall & Consult ok \\
\midrule
LLM-only & 98\% & 1.00 & 100\% & -- & -- & -- \\
SLM-only & 67\% & 0.13 & 0\% & -- & -- & -- \\
\textbf{Joule (full)} & \textbf{98\%} & \textbf{0.24} & 38\% & 53\% & 77\% & 83\% \\
No consultations & 98\% & 0.40 & 88\% & 30\% & 100\% & -- \\
Prose-only consultations & 98\% & 0.42 & 40\% & 55\% & 85\% & 40\% \\
No static checks & 96\% & 0.26 & 38\% & 53\% & 77\% & 79\% \\
Self-reported confidence & 96\% & 0.28 & 46\% & 52\% & 92\% & 74\% \\
No verification & 80\% & 0.21 & 6\% & 67\% & 15\% & -- \\
\bottomrule
\end{tabular}
\end{table}

\paragraph{Consultations on long tasks.}
On 30 further four-problem bundles (the small model alone solves 12\%), returning the edit raises the share of consultations after which the small model finishes from 13\% to 32\%, reduces handoffs from 20 to 14, and improves success from 60\% to 67\% at 22\% lower cost.
Handoffs still carry most of the load on long tasks; consultations now carry a meaningful part of it.

\paragraph{A 2026 model lineup.}
The comparison holds for current models.
With Qwen3.5~9B as the small model, GPT-5.6~Luna (minimal reasoning) as the middle rung, and Claude~Sonnet~5 on top, on 40 further MBPP problems every escalating strategy reaches Sonnet's 100\%; the ladder does so at 8\% of Sonnet's cost, the two-tier policy at 15\%, and the FrugalGPT-style cascade at 54\%.
The ladder used the middle rung on 6 problems and Sonnet on 3, with precision 78\% and recall 100\%.
The 9B model alone solves 79\% of the slice (Llama~3.1~8B: 62\%); the middle model alone, at minimal reasoning effort, only 60\%, yet as a rung above the 9B model it was enough to keep Sonnet out of most tasks.
On a cheaper stack (Qwen3.5~9B $
ightarrow$ DeepSeek~V4~Flash $
ightarrow$ DeepSeek~V4~Pro, 100 further problems, three small-model-only runs per problem), every escalating strategy matches the top model's 98\%; over three seeds the two-tier policy reaches 98\% $\pm$ 1.5 at 37\% of the top model's cost and the ladder 96\% $\pm$ 2.0 at 27\% $\pm$ 5, against 59\% for the FrugalGPT-style cascade.
The middle model alone reaches 97\% for 13\%, cheaper than the small model itself because it finishes in fewer steps; on problems this small a capable cheap model needs no escalation, and the small-model-first design pays for the small model's turns before reaching it.

\paragraph{Real repositories.}
To test the policy on repository work rather than single functions, we run SWE-bench Lite~\cite{jimenez2024swebench} instances in their official evaluation images: the agent sees the issue text and edits, reads, and runs commands inside the container; the hidden tests are applied afterwards and an instance counts as resolved only if every FAIL\_TO\_PASS and PASS\_TO\_PASS test passes.
On 15 instances (Django, pytest, pylint) with Qwen3.5~9B $\rightarrow$ Gemini~2.5~Flash $\rightarrow$ Gemini~2.5~Pro and a \$0.20 cost ceiling per task, the small model alone resolves 2, the middle model alone 2, and the ladder 7, four of which neither model resolved on its own; three were resolved by the 9B model without escalating, at one to three cents each.
A first run of the same slice with Claude~Sonnet~5 on top and the earlier rule that sends reasoning breakdowns straight to the top rung resolved 8 of 15 at 2.4$\times$ the cost, because the 9B model's unparseable first turns on 8 instances went directly to the frontier model; climbing one rung instead is now the default.
That run also exposed defects that function-level problems never trigger: search commands that exit 1 were counted as failed steps, whole-file rewrites of 600-line files exceeded the output cap, and a run that finished without changing any file was accepted.
The remaining weakness is the opposite of over-escalation: a small model that reads files for thirty steps without ever failing produces no evidence and exhausts its step budget.
We have since added an exploration-stall rule (eight consecutive successful steps that neither change a file nor verify anything trigger a consult); it is implemented and unit-tested but not yet measured on the repository slice.
Fifteen instances is a small sample; the middle model alone moved between 6 and 2 resolved across the two runs with no change that affects it.

\paragraph{Small-model compatibility.}
Nine small models from 3B to 70B drove the same step agent on 20 further problems with Gemini~2.5~Flash as the large model.
All reach 90--100\% success under adaptive escalation; the 7B--8B open models (Qwen~2.5~7B, Llama~3.1~8B) finish 100\% at about \$0.003 per problem, a third of the large model's cost, with essentially no unparseable turns.

\paragraph{Trigger ablation.}
On a 30-problem development slice, allowing the small model one verification retry before a consult, and treating an improving test pass fraction as progress, reduced cost per task by 35\% at unchanged success and raised precision from 44\% to 75\% against the same labels, at the cost of recall (100\% to 82\%).

\paragraph{Can the trigger be learned?}
We trained a logistic model over the confidence signals (plus step index, failure count, and verification-failure streak) on every small-model-only trajectory from the runs above, 1{,}092 trajectories over 364 tasks, using the run outcome as the label for each decision point and holding out 30\% of tasks.
The signals are informative: held-out AUC for predicting eventual small-model success is 0.89 per decision point.
They are not informative early: AUC is 0.51 at the first decision and 0.89 at the second, so essentially all of the information arrives with the first tool result and verifier outcome, which is exactly what the rules act on.
Replayed as a trigger (escalate when predicted success falls below $\theta$) on the held-out trajectories, the learned model reproduces the rule-based policy at its best threshold (precision 68\% and recall 91\% versus 69\% and 90\%) rather than improving on it.
We read this as evidence that the rule set is already close to what the current signals support; better triggers need richer evidence at the step, such as the content of failing tests, not a better classifier over the same signals.

\paragraph{Limitations.}
Three coding workloads, one small model in the main comparison, a 15-instance repository slice, and rule-based triggers whose thresholds were set by hand.
The three-rung and repository results are single-seed.
Latency is not a win: adaptive runs take 1.7$\times$ longer than the large model alone, because the small model takes more turns and each is slower.
The counterfactual labels are themselves estimates from three runs of a stochastic model.
